% Makra obrazku kapitoly 12 (Prostorova slozitost) -- opsano z ../skripta/kapitoly/12-prostorova-slozitost.tex
\newcommand{\savrow}[5]{% #1 posun, #2 popisek, #3 zkouseny mid, #4 radek patra 0, #5 radek patra 1
  \begin{scope}[yshift=#1]
    \node[font=\scriptsize\bfseries,anchor=west] at (-0.4,0.75) {#2};
    \foreach \i in {1,...,4} \node[cf] (c\i) at ({(\i-1)*1.3},0) {$C_\i$};
    \foreach \a/\b in {1/2,2/3,3/4} \draw[->,thick,black!60] (c\a)--(c\b);
    \node[cf,fill=hl] at (c#3) {$C_#3$};
    \node[anchor=west] at (5.2,-0.05) {\scriptsize
      \begin{tabular}{|c|c|c|c|}\hline patro & konce & $t$ & $C_{\mathrm{mid}}$\\\hline
      0 & #4\\\hline 1 & #5\\\hline\end{tabular}};
  \end{scope}}
\newcommand{\ndtree}[5]{% #1 posun, #2-#5 listy (ok/ko)
  \begin{scope}[xshift=#1]
    \node[lf,fill=white] (r) at (1.5,1.6) {};
    \node[lf,fill=white] (p) at (0.75,0.8) {}; \node[lf,fill=white] (q) at (2.25,0.8) {};
    \node[#2] (l1) at (0.3,0) {}; \node[#3] (l2) at (1.2,0) {};
    \node[#4] (l3) at (1.8,0) {}; \node[#5] (l4) at (2.7,0) {};
    \draw[thick,black!60] (r)--(p) (r)--(q) (p)--(l1) (p)--(l2) (q)--(l3) (q)--(l4);
  \end{scope}}
